\documentclass[11pt]{article}
\usepackage{mystyle}

\title{Rosen --- \S 2.4: Sequences and Summations\\ \large Selected Exercises}
\author{Alston}
\date{Semester 2, 2026}

\begin{document}
\maketitle

% ======================================================================
\begin{exercise}{11}
    Let $a_n = 2^n + 5 \cdot 3^n$ for $n = 0, 1, 2, \dots$.
    \begin{enumerate}[label=(\alph*)]
        \item Find $a_0$, $a_1$, $a_2$, $a_3$, and $a_4$.
        \item Show that $a_2 = 5a_1 - 6a_0$, $a_3 = 5a_2 - 6a_1$, and
        $a_4 = 5a_3 - 6a_2$.
        \item Show that $a_n = 5a_{n-1} - 6a_{n-2}$ for all integers $n$
        with $n \geq 2$.
    \end{enumerate}
\end{exercise}

% ======================================================================
\begin{exercise}{27}
    Show that if $a_n$ denotes the $n$th positive integer that is not a
    perfect square, then $a_n = n + \{\sqrt{n}\}$, where $\{x\}$ denotes
    the integer closest to the real number $x$.
\end{exercise}

% ======================================================================
\begin{exercise}{35}
    Show that
    \[
    \sum_{j=1}^{n} (a_j - a_{j-1}) = a_n - a_0,
    \]
    where $a_0, a_1, \dots, a_n$ is a sequence of real numbers. This
    type of sum is called \emph{telescoping}.
\end{exercise}

% ======================================================================
\begin{exercise}{37}
    Sum both sides of the identity $k^2 - (k-1)^2 = 2k - 1$ from $k = 1$
    to $k = n$ and use Exercise 35 to find
    \begin{enumerate}[label=(\alph*)]
        \item a formula for $\displaystyle\sum_{k=1}^{n} (2k-1)$ (the
        sum of the first $n$ odd natural numbers).
        \item a formula for $\displaystyle\sum_{k=1}^{n} k$.
    \end{enumerate}
\end{exercise}

% ======================================================================
\begin{exercise}{38}
    Use the technique given in Exercise 35, together with the result of
    Exercise 37(b), to derive the formula for
    $\displaystyle\sum_{k=1}^{n} k^2$ given in Table 2.
    [\emph{Hint}: Take $a_k = k^3$ in the telescoping sum in Exercise
    35.]
\end{exercise}

% ======================================================================
\begin{exercise}{41}
    Find a formula for $\displaystyle\sum_{k=0}^{m} \lfloor\sqrt{k}\rfloor$,
    when $m$ is a positive integer.
\end{exercise}

% ======================================================================
\begin{exercise}{42}
    Find a formula for $\displaystyle\sum_{k=0}^{m} \lfloor\sqrt[3]{k}\rfloor$,
    when $m$ is a positive integer.
\end{exercise}

\end{document}